\documentclass[10pt,a4paper]{amsart}
\usepackage[margin=19mm]{geometry}
\usepackage{amsmath,amssymb,mathtools}
\usepackage[scr=rsfs,cal=euler]{mathalfa}
\usepackage{xfrac}
\usepackage[unicode,hidelinks,bookmarks=false]{hyperref}
\newcommand{\ie}{i.e.\ }
\newcommand{\cf}{cf.\ }
\newcommand{\resp}{respectively\ }
\newcommand{\Subsection}{\subsection}
\providecommand{\nobreakdash}{\nobreak\hbox{-}}
\setlength{\emergencystretch}{2em}
\multlinegap0pt
\usepackage{amssymb}
\usepackage{xcolor}
\usepackage[shortlabels]{enumitem}
\usepackage{tikz}
\newtheorem{proposition}{Proposition}
\newtheorem{lemma}{Lemma}
\usepackage{cleveref}
\crefname{proposition}{Proposition}{Propositions}
\crefname{lemma}{Lemma}{Lemmas}
\newcommand{\N}{\mathbb{N}}
\newcommand{\R}{\mathbb{R}}
\newcommand{\Tr}{\operatorname{Tr}}
\newcommand{\Vol}{\operatorname{Vol}}
\renewcommand{\d}{\operatorname{d}\!}
\newcommand{\del}{\partial}
\title{A stabilization result for $\mathbb{Z}_2$-harmonic 1-forms by constructing solutions on closed 3-manifolds with long cylindrical necks}
\author{Willem Adriaan Salm}
\date{Source: arXiv:2410.07015v2 (CC BY 4.0)}
\begin{document}
\maketitle

\noindent\emph{Source and licence.} A stabilization result for $\mathbb{Z}_2$-harmonic 1-forms by constructing solutions on closed 3-manifolds with long cylindrical necks. Version arXiv:2410.07015v2 (CC BY 4.0), distributed by arXiv under the Creative Commons Attribution 4.0 International licence, http://creativecommons.org/licenses/by/4.0/, as recorded in the arXiv licence metadata for this exact version and shown on the abstract page https://arxiv.org/abs/2410.07015v2. Full licence notice: \texttt{licenses/arxiv-2410.07015-CC-BY-4.0.txt}.\par\smallskip

\noindent\textit{Editorial scope.} Counted unit: the article's lemma lem:proof-assumptions-non-connected-case (source0.tex lines 1567--1581). The article's complete original proof of it is delivered with it (source0.tex lines 1582--1638, one \texttt{proof} environment of 57 lines). No mathematics is written, repaired or re-derived: the statement and the proof are the article's own lines, in the article's own notation, and no retained line is substituted or re-flowed.  Retained uncounted as the source-local context the proof needs: lines 1384--1391; lines 1472--1480; lines 1525--1534.  Nothing was omitted from any retained span, so the package carries no omission record.  No retained line contains a citation, so the wrapper carries no bibliography and no bibliography entry.  No retained line contains a figure environment, an image-inclusion command or drawing code, so no image is delivered and the package declares no asset dependency.  Editorial formatting only, in addition: an AMS \texttt{amsart} preamble that restates the article's own declarations and macros used by the retained lines, namely the environments \texttt{proposition}, \texttt{lemma} on separate counters, together with the standard packages those lines need.  The retained environments print in this document as Proposition 1, Lemma 1 in order of appearance, whereas the article prints the same items with the numbering of its own full document.  The complete native source of the article as distributed in the arXiv e-print for this exact version is bundled unchanged as \texttt{sources/166566/source0.tex}.
\begin{proposition}
	\label{prop:variation-metric-v-cyl}
	For any fixed $\omega$ with a non-trivial cohomology class and any positive odd number $n$, one can always perturb the metric on the neck near the interior region such that 
	$$
	\Re(v^{\mathrm{cyl}}_{n0}(\omega, \Sigma_i)) \not = 0 \quad\text{and/or}\quad \Im(v^{\mathrm{cyl}}_{n0}(\omega, \Sigma_i)) \not = 0
	$$
	for any connected component $\Sigma_i$ of $\Sigma$. 
\end{proposition}

\begin{lemma}
	\label{lem:variation:lemma2.2-he}
	Let $U$ be an open neighbourhood in $\hat{M}$, disjoint from the support of $\d^* \omega$. Suppose that for all $T$ supported on $U$,
	\begin{equation}
		\label{eq:variation:lemma2.2-he}
		  \int_{\hat{M}} (\Tr(TS) - \frac{1}{2}\Tr(T)\Tr(S))\Vol^{g_{\mathrm{cyl}}} = 0.
	\end{equation}
	Then, wherever $\d G \not= 0$ on $U$, we have $\d v^{\mathrm{cyl}} = 0$.
\end{lemma}

\begin{lemma}
	\label{lem:proof-assumptions-connected-case}
	Assume $\Sigma$ is connected.
	Let $E \subset H^1_-(\hat{M})$ be a $2$ dimensional subspace. For a generic metric, there is a basis $\{ [\omega^\Re], [\omega^\Im]\}$ of $E$ such that 
	$$
	v^{\mathrm{cyl}}_{10}(\omega^\Re, \Sigma) = 1
	\quad \text{and} \quad
	v^{\mathrm{cyl}}_{10}(\omega^\Im, \Sigma) = i
	$$
\end{lemma}

\begin{lemma}
	\label{lem:proof-assumptions-non-connected-case}
	Let $p$ be the number of connected components of $\Sigma$ and let $E \subset H^1_-(\hat{M})$ be a $2p$-dimensional subspace. For a generic metric, there is a basis $\{ [\omega_1^\Re], [\omega_1^\Im], \ldots, [\omega_p^\Re], [\omega_p^\Im]\}$ of $E$ such that 
	$$
	v^{\mathrm{cyl}}_{10}(\omega_k^\Re, \Sigma_l) = \begin{cases}
		1 & k = l, \\
		0 & k \not= l,
	\end{cases}
	\quad \text{and} \quad
	v^{\mathrm{cyl}}_{10}(\omega_k^\Im, \Sigma_l) = \begin{cases}
		i & k = l, \\
		0 & k \not= l.
	\end{cases}
	$$
\end{lemma}

\begin{proof}
	In \cref{lem:proof-assumptions-connected-case} we already considered the case where $\Sigma$ is connected. Now consider the case where $\Sigma$ is not connected. Let $\Sigma_1, \ldots, \Sigma_p$ be the path-connected components of $\Sigma$.
	Consider the map $V \colon E \to \R^{2p}$, given by
	$$
	V([\omega]) = \begin{pmatrix}
		\Re(v^{\mathrm{cyl}}_{10}(\omega, \Sigma_1))\\ \Im(v^{\mathrm{cyl}}_{10}(\omega, \Sigma_1)) \\
		\vdots \\
		\Re(v^{\mathrm{cyl}}_{10}(\omega, \Sigma_p))\\ \Im(v^{\mathrm{cyl}}_{10}(\omega, \Sigma_p)) 
	\end{pmatrix}.$$
	If $\{[\omega_1], \ldots [\omega_{2p}]\}$ is a basis of $E$, then $V$ can be written as
	\begin{equation}
		\label{eq:variation:matrix-V}
		V = \begin{pmatrix}
			\Re(v^{\mathrm{cyl}}_{10}(\omega_1, \Sigma_1)) 
			& \Re(v^{\mathrm{cyl}}_{10}(\omega_2, \Sigma_1)) 
			& \Re(v^{\mathrm{cyl}}_{10}(\omega_3, \Sigma_1)) 
			& \ldots \\
			\Im(v^{\mathrm{cyl}}_{10}(\omega_1, \Sigma_1)) 
			& \Im(v^{\mathrm{cyl}}_{10}(\omega_2, \Sigma_1))
			& \Im(v^{\mathrm{cyl}}_{10}(\omega_3, \Sigma_1))
			& \ldots \\
			\Re(v^{\mathrm{cyl}}_{10}(\omega_1, \Sigma_2)) 
			& \Re(v^{\mathrm{cyl}}_{10}(\omega_2, \Sigma_2)) 
			& \Re(v^{\mathrm{cyl}}_{10}(\omega_3, \Sigma_2)) 
			& \ldots \\
			\vdots & \vdots & \vdots & \ddots
		\end{pmatrix}.
	\end{equation}
	Again, if $\det (V) \not = 0$, then we can find a basis of $E$ that makes $V$ the identity matrix, proving the lemma. So assume that $\det V = 0$ and the lemma is false.
	
	Let $V^k$ be a square $k$ by $k$ matrix that is constructed by restricting $V$ to the first $k$ rows and columns. Consider the first instance of $k \in \N$ such that $\det(V^{k-1}) \not = 0$ while $\det V^k = 0$.
	We can write $V^k$ as a block matrix
	$$
	V^k = \begin{pmatrix}
		V^{k-1} & U \\
		W & x
	\end{pmatrix}
	$$ 
	and using some basis transformations, we can always assume that 
	$$
	V^k = \begin{pmatrix}
		\operatorname{Id} & 0 \\
		W & x
	\end{pmatrix}.
	$$ 
	By the assumption that $\det V^k = 0$, the value of $x$ has to vanish. Using the Schur complement, the determinant of $V^k$ can be written as
	$$
	\det(V^k) = \det(V^{k-1}) \cdot (x - W (V^{k-1})^{-1} U).
	$$
	So the variation of $\det(V^k)$ under the variation of the metric simplifies to
	$$
	\left.\frac{\del }{\del t}\right|_{t=0} \det(V^k) = \left.\frac{\del}{\del t}\right|_{t=0} x - W\cdot \left.\frac{\del }{\del t}\right|_{t=0} U.
	$$

	Now look at \cref{eq:variation:matrix-V}: The element $x$ and the row vector $U$ only depend on $\omega_k$, while $W$ only depends on $\omega_1, \ldots, \omega_{k-1}$. Like in \cref{lem:proof-assumptions-connected-case}, we can treat $W$ as a constant column vector and write $U$ and $x$ in terms of Poisson maps that act on $\d^* \omega_k$. The resulting Poisson map is not constant and so using \cref{lem:variation:lemma2.2-he} and repeating the proof of \cref{prop:variation-metric-v-cyl} we can reach a contradiction. This concludes that for a generic metric $V^k$ is invertible. 
	By induction, $V$ is invertible for a generic metric, which finishes the proof.
\end{proof}

\end{document}
